\documentclass[11pt]{article}

\author{Math 1030}
\date{Due Friday, April 10 by 5pm} 
\title{Homework 8}

\usepackage{graphicx,xypic}
\usepackage{amsthm}
\usepackage{amsmath,amssymb}
\usepackage{amsfonts}
\usepackage{xcolor}
\usepackage[margin=1in]{geometry}
\usepackage[shortlabels]{enumitem}
\newtheorem{problem}{Problem}
\renewcommand*{\proofname}{{\color{blue}Solution}}


\setlength{\parindent}{0pt}
\setlength{\parskip}{1.25ex}


\begin{document}

\maketitle

% DON'T LEAVE THIS BLANK! 
{\bf\Large Your Name:} 

% DON'T FORGET YOUR COLLABORATORS! 
Collaborator names: 

\vspace{.3in}
Topics covered: Ramsey theory, pigeonhole principle

Instructions {\bf (read these carefully!)}: 
\begin{itemize}
\item This assignment must be submitted on Gradescope by the due date. 
\item If you collaborate with other students (this is encouraged),  list your collaborators above. 
\item Homework is graded anonymously, so avoid putting your name on each page of the assignment.
\item If you are stuck, please ask for help (from me, Rafi, or a classmate). Use Ed discussions!  
\item You may freely use any fact proved in class (mention it in the appropriate spot). You \emph{may not} freely use facts from the book or other sources.
\item As a general rule, try to restrict your solution to each problem to a single page. Often solutions can be short, while still being complete. If your solution is very long, think more about how you might express it more concisely.
\end{itemize}
\pagebreak 


\begin{problem}
Use the pigeonhole principle to prove the following statements. 
\begin{enumerate}[(a)]
\item Every set of $n$ integers $\{a_1,\ldots,a_n\}$ contains a nonempty subset whose sum is divisible by $n$. \footnote{Hint: consider the partial sums $S_i=a_1+\cdots+a_i$.}\footnote{Remark: The set $\{1,n+1,2n+1,\ldots,(n-1)n+1\}$ shows you cannot improve this result to a set of $n-1$ integers.}
\item Given $x\in\mathbb R$, prove that at least one of $\{x,2x,\ldots,(n-1)x\}$ differs by at most $1/n$ from an integer. \footnote{Hint: consider the fractional parts of these integers.}
\end{enumerate} 
\end{problem}

\begin{proof}

\end{proof}

\pagebreak 

\begin{problem}
Each of two concentric discs has $20$ radial sections of equal size. For each disc, $10$ sections are painted red and $10$ blue, in some arrangement. Prove that the two discs can be aligned so that at least $10$ sections on the inner disc match colors with the corresponding sections on the outer disc. \footnote{Hint: There is a very short solution using the pigeonhole principle. Fixing a section of the outer disk, how many positions of the inner disk give a matching in this one ``coordinate"?}
\end{problem}

\begin{proof}

\end{proof}

\pagebreak 

\begin{problem}
Fix $r\ge2$ and let $R(k,\ldots,k)$ (with $r$ copies of $k$) denote the minimal $n$ so that any $r$-coloring of the edges of $K_n$ contains at a monochromatic $K_k$ (i.e.\ we are generalizing the Ramsey numbers to more colors). Prove that every $r$ coloring of $1,\ldots,R(3,\ldots,3)$ (with $r$ copies of $3$) contains a monochromatic $x,y,z$ so that $x+y=z$. \footnote{Hint: set $n=R(3,\ldots,3)$. Given a coloring of $1,\ldots,n$, construct an edge 2-coloring of $K_n$ so that a monochromatic triangle in $K_n$ corresponds to monochromatic $x,y,z$ with $x+y=z$.} 
\end{problem}

\begin{proof}

\end{proof}

\pagebreak 

\begin{problem}
A theorem similar to Negami's theorem says that for any embedding of $K_6$ in $\mathbb R^3$, there exists a pair of triangles that form a Hopf link.\footnote{Google this.} Verify this result for the following set of 6 points in $\mathbb R^3$. 
\[A=(0,3,2), \>\>\>B=(-2,-6,0), \>\>\>C=(6,3,2), \>\>\>D=(-6,-10,7), \>\>\>E=(-9,-6,9),\>\>\>F=(3,-1,9)\]
To help solve this, use this Geo-gebra notebook: 
\texttt{https://www.geogebra.org/calculator/rrbafuh6} \footnote{Click the circles on the left panel to make line segments appear/disappear.} Include a screenshot in your solution. 
\end{problem}

\begin{proof}

\end{proof}

\pagebreak 

\begin{problem}[Bonus]
Give a set of eight points $S$ in the plane, no three on a line, with the property that $S$ does not contain the vertex set of any convex $5$-gon. Be sure to give an explanation.\footnote{In fact, this is the best you can do: any collection of $9$ points in the plane (no three collinear) contains the vertices of a convex pentagon.} 
\end{problem}

\begin{proof}

\end{proof}




\end{document}